% Options for packages loaded elsewhere
\PassOptionsToPackage{unicode}{hyperref}
\PassOptionsToPackage{hyphens}{url}
\documentclass[
]{article}
\usepackage{xcolor}
\usepackage[margin=1in]{geometry}
\usepackage{amsmath,amssymb}
\setcounter{secnumdepth}{-\maxdimen} % remove section numbering
\usepackage{iftex}
\ifPDFTeX
  \usepackage[T1]{fontenc}
  \usepackage[utf8]{inputenc}
  \usepackage{textcomp} % provide euro and other symbols
\else % if luatex or xetex
  \usepackage{unicode-math} % this also loads fontspec
  \defaultfontfeatures{Scale=MatchLowercase}
  \defaultfontfeatures[\rmfamily]{Ligatures=TeX,Scale=1}
\fi
\usepackage{lmodern}
\ifPDFTeX\else
  % xetex/luatex font selection
\fi
% Use upquote if available, for straight quotes in verbatim environments
\IfFileExists{upquote.sty}{\usepackage{upquote}}{}
\IfFileExists{microtype.sty}{% use microtype if available
  \usepackage[]{microtype}
  \UseMicrotypeSet[protrusion]{basicmath} % disable protrusion for tt fonts
}{}
\makeatletter
\@ifundefined{KOMAClassName}{% if non-KOMA class
  \IfFileExists{parskip.sty}{%
    \usepackage{parskip}
  }{% else
    \setlength{\parindent}{0pt}
    \setlength{\parskip}{6pt plus 2pt minus 1pt}}
}{% if KOMA class
  \KOMAoptions{parskip=half}}
\makeatother
\usepackage{color}
\usepackage{fancyvrb}
\newcommand{\VerbBar}{|}
\newcommand{\VERB}{\Verb[commandchars=\\\{\}]}
\DefineVerbatimEnvironment{Highlighting}{Verbatim}{commandchars=\\\{\}}
% Add ',fontsize=\small' for more characters per line
\usepackage{framed}
\definecolor{shadecolor}{RGB}{248,248,248}
\newenvironment{Shaded}{\begin{snugshade}}{\end{snugshade}}
\newcommand{\AlertTok}[1]{\textcolor[rgb]{0.94,0.16,0.16}{#1}}
\newcommand{\AnnotationTok}[1]{\textcolor[rgb]{0.56,0.35,0.01}{\textbf{\textit{#1}}}}
\newcommand{\AttributeTok}[1]{\textcolor[rgb]{0.13,0.29,0.53}{#1}}
\newcommand{\BaseNTok}[1]{\textcolor[rgb]{0.00,0.00,0.81}{#1}}
\newcommand{\BuiltInTok}[1]{#1}
\newcommand{\CharTok}[1]{\textcolor[rgb]{0.31,0.60,0.02}{#1}}
\newcommand{\CommentTok}[1]{\textcolor[rgb]{0.56,0.35,0.01}{\textit{#1}}}
\newcommand{\CommentVarTok}[1]{\textcolor[rgb]{0.56,0.35,0.01}{\textbf{\textit{#1}}}}
\newcommand{\ConstantTok}[1]{\textcolor[rgb]{0.56,0.35,0.01}{#1}}
\newcommand{\ControlFlowTok}[1]{\textcolor[rgb]{0.13,0.29,0.53}{\textbf{#1}}}
\newcommand{\DataTypeTok}[1]{\textcolor[rgb]{0.13,0.29,0.53}{#1}}
\newcommand{\DecValTok}[1]{\textcolor[rgb]{0.00,0.00,0.81}{#1}}
\newcommand{\DocumentationTok}[1]{\textcolor[rgb]{0.56,0.35,0.01}{\textbf{\textit{#1}}}}
\newcommand{\ErrorTok}[1]{\textcolor[rgb]{0.64,0.00,0.00}{\textbf{#1}}}
\newcommand{\ExtensionTok}[1]{#1}
\newcommand{\FloatTok}[1]{\textcolor[rgb]{0.00,0.00,0.81}{#1}}
\newcommand{\FunctionTok}[1]{\textcolor[rgb]{0.13,0.29,0.53}{\textbf{#1}}}
\newcommand{\ImportTok}[1]{#1}
\newcommand{\InformationTok}[1]{\textcolor[rgb]{0.56,0.35,0.01}{\textbf{\textit{#1}}}}
\newcommand{\KeywordTok}[1]{\textcolor[rgb]{0.13,0.29,0.53}{\textbf{#1}}}
\newcommand{\NormalTok}[1]{#1}
\newcommand{\OperatorTok}[1]{\textcolor[rgb]{0.81,0.36,0.00}{\textbf{#1}}}
\newcommand{\OtherTok}[1]{\textcolor[rgb]{0.56,0.35,0.01}{#1}}
\newcommand{\PreprocessorTok}[1]{\textcolor[rgb]{0.56,0.35,0.01}{\textit{#1}}}
\newcommand{\RegionMarkerTok}[1]{#1}
\newcommand{\SpecialCharTok}[1]{\textcolor[rgb]{0.81,0.36,0.00}{\textbf{#1}}}
\newcommand{\SpecialStringTok}[1]{\textcolor[rgb]{0.31,0.60,0.02}{#1}}
\newcommand{\StringTok}[1]{\textcolor[rgb]{0.31,0.60,0.02}{#1}}
\newcommand{\VariableTok}[1]{\textcolor[rgb]{0.00,0.00,0.00}{#1}}
\newcommand{\VerbatimStringTok}[1]{\textcolor[rgb]{0.31,0.60,0.02}{#1}}
\newcommand{\WarningTok}[1]{\textcolor[rgb]{0.56,0.35,0.01}{\textbf{\textit{#1}}}}
\usepackage{graphicx}
\makeatletter
\newsavebox\pandoc@box
\newcommand*\pandocbounded[1]{% scales image to fit in text height/width
  \sbox\pandoc@box{#1}%
  \Gscale@div\@tempa{\textheight}{\dimexpr\ht\pandoc@box+\dp\pandoc@box\relax}%
  \Gscale@div\@tempb{\linewidth}{\wd\pandoc@box}%
  \ifdim\@tempb\p@<\@tempa\p@\let\@tempa\@tempb\fi% select the smaller of both
  \ifdim\@tempa\p@<\p@\scalebox{\@tempa}{\usebox\pandoc@box}%
  \else\usebox{\pandoc@box}%
  \fi%
}
% Set default figure placement to htbp
\def\fps@figure{htbp}
\makeatother
\setlength{\emergencystretch}{3em} % prevent overfull lines
\providecommand{\tightlist}{%
  \setlength{\itemsep}{0pt}\setlength{\parskip}{0pt}}
\usepackage{booktabs}
\usepackage{longtable}
\usepackage{array}
\usepackage{multirow}
\usepackage{wrapfig}
\usepackage{float}
\usepackage{colortbl}
\usepackage{pdflscape}
\usepackage{tabu}
\usepackage{threeparttable}
\usepackage{threeparttablex}
\usepackage[normalem]{ulem}
\usepackage{makecell}
\usepackage{xcolor}
\usepackage{bookmark}
\IfFileExists{xurl.sty}{\usepackage{xurl}}{} % add URL line breaks if available
\urlstyle{same}
\hypersetup{
  pdftitle={Case Study 2},
  pdfauthor={First Name, Last Name},
  hidelinks,
  pdfcreator={LaTeX via pandoc}}

\title{Case Study 2}
\author{First Name, Last Name}
\date{2026-04-02}

\begin{document}
\maketitle

\emph{The .Rmd} \textbf{and} \emph{.html (or .pdf) should be uploaded in
TUWEL by the deadline. Refrain from using explanatory comments in the R
code chunks but write them as text instead. Points will be deducted if
the submitted file is not in a decent form.}

\textbf{DISCLAIMER}: In case students did not contribute equally,
include a disclaimer stating what each student's contribution was.

The CIA World Factbook provides intelligence on various aspects of 266
world entities, including history, people, government, economy, energy,
geography, environment, communications, transportation, military,
terrorism, and transnational issues. This case study involves analyzing
world data from 2020, focusing on:

\begin{itemize}
\tightlist
\item
  \textbf{Education Expenditure (\% of GDP)}
\item
  \textbf{Youth Unemployment Rate (15-24 years)}
\item
  \textbf{Net Migration Rate} (difference between the number of people
  entering and leaving a country per 1,000 persons)
\end{itemize}

The data was sourced from the
\href{https://www.cia.gov/the-world-factbook/about/archives/}{CIA World
Factbook Archives}. You are required to use \texttt{dplyr} for data
manipulation, while any package can be used for importing data.

\section{Tasks:}\label{tasks}

\subsection{a. Data Import and
Cleaning}\label{a.-data-import-and-cleaning}

Load the following datasets from TUWEL and ensure that missing values
are handled correctly and column names are clear. Each dataset should
ultimately contain only two columns: \textbf{country} and the respective
variable. Note that some data sets also contain information on the year
when the value was last updated.

\begin{itemize}
\item
  \texttt{rawdata\_369.txt} which contains the (estimated) public
  expenditure on education as a percent of GDP. \emph{Pay attention! The
  delimiter is 2 or more white spaces (one space would not work as it
  would separate country names which contain a space); you have to skip
  the first two lines}.
\item
  \texttt{rawdata\_373.csv} which contains the (estimated) youth
  unemployment rate (15-24) per country
\item
  \texttt{rawdata\_347.txt} which contains (estimated) net migration
  rate per country.
\end{itemize}

\begin{Shaded}
\begin{Highlighting}[]
\FunctionTok{library}\NormalTok{(dplyr)}
\end{Highlighting}
\end{Shaded}

\begin{verbatim}
## 
## Attache Paket: 'dplyr'
\end{verbatim}

\begin{verbatim}
## Die folgenden Objekte sind maskiert von 'package:stats':
## 
##     filter, lag
\end{verbatim}

\begin{verbatim}
## Die folgenden Objekte sind maskiert von 'package:base':
## 
##     intersect, setdiff, setequal, union
\end{verbatim}

\begin{Shaded}
\begin{Highlighting}[]
\FunctionTok{library}\NormalTok{(readr)}
\FunctionTok{library}\NormalTok{(stringr)}
\FunctionTok{library}\NormalTok{(tibble)}
\FunctionTok{library}\NormalTok{(tidyr)}
\FunctionTok{library}\NormalTok{(kableExtra)}
\end{Highlighting}
\end{Shaded}

\begin{verbatim}
## 
## Attache Paket: 'kableExtra'
\end{verbatim}

\begin{verbatim}
## Das folgende Objekt ist maskiert 'package:dplyr':
## 
##     group_rows
\end{verbatim}

\begin{Shaded}
\begin{Highlighting}[]
\CommentTok{\# Load rawdata\_369.txt (public expenditure on education)}
\NormalTok{education\_data }\OtherTok{\textless{}{-}} \FunctionTok{read\_lines}\NormalTok{(}\StringTok{"rawdata\_369.txt"}\NormalTok{, }\AttributeTok{skip =} \DecValTok{2}\NormalTok{)}

\NormalTok{education\_expenditure }\OtherTok{\textless{}{-}} \FunctionTok{tibble}\NormalTok{(}
  \AttributeTok{country =} \FunctionTok{substr}\NormalTok{(education\_data, }\DecValTok{9}\NormalTok{, }\DecValTok{72}\NormalTok{),}
  \AttributeTok{education\_expenditure =} \FunctionTok{as.numeric}\NormalTok{(}\FunctionTok{substr}\NormalTok{(education\_data, }\DecValTok{73}\NormalTok{, }\DecValTok{84}\NormalTok{))}
\NormalTok{) }\SpecialCharTok{\%\textgreater{}\%}
  \FunctionTok{mutate}\NormalTok{(}\AttributeTok{country =} \FunctionTok{str\_trim}\NormalTok{(country))}

\FunctionTok{glimpse}\NormalTok{(education\_expenditure)}
\end{Highlighting}
\end{Shaded}

\begin{verbatim}
## Rows: 170
## Columns: 2
## $ country               <chr> "Cuba", "Micronesia, Federated States of", "Solo~
## $ education_expenditure <dbl> 12.8, 12.5, 9.9, 8.8, 7.9, 7.8, 7.7, 7.7, 7.6, 7~
\end{verbatim}

\begin{Shaded}
\begin{Highlighting}[]
\CommentTok{\# Load rawdata\_373.csv (youth unemployment rate)}
\NormalTok{youth\_unemployment\_rate }\OtherTok{\textless{}{-}} \FunctionTok{read\_csv}\NormalTok{(}\StringTok{"rawdata\_373.csv"}\NormalTok{, }\AttributeTok{show\_col\_types =} \ConstantTok{FALSE}\NormalTok{) }\SpecialCharTok{\%\textgreater{}\%}
  \FunctionTok{rename}\NormalTok{(}
    \AttributeTok{country =}\NormalTok{ country\_name,}
    \AttributeTok{youth\_unemployment\_rate =}\NormalTok{ youth\_unempl\_rate}
\NormalTok{  ) }\SpecialCharTok{\%\textgreater{}\%}
  \FunctionTok{mutate}\NormalTok{(}\AttributeTok{country =} \FunctionTok{str\_trim}\NormalTok{(country))}
\FunctionTok{glimpse}\NormalTok{(youth\_unemployment\_rate)}
\end{Highlighting}
\end{Shaded}

\begin{verbatim}
## Rows: 181
## Columns: 2
## $ country                 <chr> "French Polynesia", "Kosovo", "South Africa", ~
## $ youth_unemployment_rate <dbl> 56.7, 55.4, 53.4, 48.7, 47.1, 46.2, 45.4, 42.2~
\end{verbatim}

\begin{Shaded}
\begin{Highlighting}[]
\CommentTok{\# Load rawdata\_347.txtv (net migration rate)}
\NormalTok{migration\_data }\OtherTok{\textless{}{-}} \FunctionTok{read\_lines}\NormalTok{(}\StringTok{"rawdata\_347.txt"}\NormalTok{, }\AttributeTok{skip =} \DecValTok{2}\NormalTok{)}

\NormalTok{net\_migration\_rate }\OtherTok{\textless{}{-}} \FunctionTok{tibble}\NormalTok{(}
  \AttributeTok{country =} \FunctionTok{str\_trim}\NormalTok{(}\FunctionTok{substr}\NormalTok{(migration\_data, }\DecValTok{9}\NormalTok{, }\DecValTok{72}\NormalTok{)),}
  \AttributeTok{net\_migration\_rate =} \FunctionTok{as.numeric}\NormalTok{(}\FunctionTok{substr}\NormalTok{(migration\_data, }\DecValTok{73}\NormalTok{, }\DecValTok{84}\NormalTok{))}
\NormalTok{) }\SpecialCharTok{\%\textgreater{}\%}
  \FunctionTok{mutate}\NormalTok{(}\AttributeTok{country =} \FunctionTok{str\_trim}\NormalTok{(country))}
\FunctionTok{glimpse}\NormalTok{(net\_migration\_rate)}
\end{Highlighting}
\end{Shaded}

\begin{verbatim}
## Rows: 227
## Columns: 2
## $ country            <chr> "Syria", "British Virgin Islands", "Luxembourg", "C~
## $ net_migration_rate <dbl> 27.1, 15.5, 13.3, 13.0, 11.8, 11.1, 10.6, 8.9, 8.4,~
\end{verbatim}

\textbf{Conclusion:}

The files rawdata\_369.txt and rawdata\_347 contain a lot of white
spaces between the columns. In order to correctly import the data, we
can not blindly use the read\_table function, but we first have to call
read\_lines and then substring the appropriate data from the lines into
the fitting column and finally trim country in order to get rid of the
tailing white spaces.

\subsection{b. Merging Raw Data}\label{b.-merging-raw-data}

Merge the datasets using \texttt{dplyr} on a unique key and retain the
union of all observations.

\begin{itemize}
\tightlist
\item
  What key are you using for merging?
\item
  Return the dimensions of the merged dataset.
\end{itemize}

\begin{Shaded}
\begin{Highlighting}[]
\NormalTok{country\_data }\OtherTok{\textless{}{-}}\NormalTok{ education\_expenditure }\SpecialCharTok{\%\textgreater{}\%}
  \FunctionTok{full\_join}\NormalTok{(youth\_unemployment\_rate, }\AttributeTok{by =} \StringTok{"country"}\NormalTok{) }\SpecialCharTok{\%\textgreater{}\%}
  \FunctionTok{full\_join}\NormalTok{(net\_migration\_rate, }\AttributeTok{by =} \StringTok{"country"}\NormalTok{)}
\FunctionTok{dim}\NormalTok{(country\_data)}
\end{Highlighting}
\end{Shaded}

\begin{verbatim}
## [1] 227   4
\end{verbatim}

\begin{Shaded}
\begin{Highlighting}[]
\FunctionTok{glimpse}\NormalTok{(country\_data)}
\end{Highlighting}
\end{Shaded}

\begin{verbatim}
## Rows: 227
## Columns: 4
## $ country                 <chr> "Cuba", "Micronesia, Federated States of", "So~
## $ education_expenditure   <dbl> 12.8, 12.5, 9.9, 8.8, 7.9, 7.8, 7.7, 7.7, 7.6,~
## $ youth_unemployment_rate <dbl> 6.1, 18.9, 1.3, NA, 9.7, 9.4, 6.1, 9.4, 15.3, ~
## $ net_migration_rate      <dbl> -3.7, -20.9, -1.6, 0.0, 4.0, 2.8, 3.3, -1.2, -~
\end{verbatim}

\begin{Shaded}
\begin{Highlighting}[]
\FunctionTok{head}\NormalTok{(country\_data)}
\end{Highlighting}
\end{Shaded}

\begin{verbatim}
## # A tibble: 6 x 4
##   country        education_expenditure youth_unemployment_r~1 net_migration_rate
##   <chr>                          <dbl>                  <dbl>              <dbl>
## 1 Cuba                            12.8                    6.1               -3.7
## 2 Micronesia, F~                  12.5                   18.9              -20.9
## 3 Solomon Islan~                   9.9                    1.3               -1.6
## 4 Montserrat                       8.8                   NA                  0  
## 5 Norway                           7.9                    9.7                4  
## 6 Denmark                          7.8                    9.4                2.8
## # i abbreviated name: 1: youth_unemployment_rate
\end{verbatim}

\textbf{Conclusion:}

The union of all observations indicates we need to use full\_join to
retain all observations.

What key are you using for merging?

We can use country as a unique key for merging, due to all the imported
data inside the tibbles having one row per country and the measured
variable for that country.

Return the dimensions of the merged dataset:

By using the dim() function on the resulting tibble we are returned the
dimensions 227x4 indicating 227 rows with observations and 4 columns
with variables.

\subsection{c.~Enriching Data with Income
Classification}\label{c.-enriching-data-with-income-classification}

Obtain country income classification (low, lower-middle, upper-middle,
high) from the
\href{https://datahelpdesk.worldbank.org/knowledgebase/articles/906519}{World
Bank} and merge it with the dataset.

\begin{itemize}
\item
  Identify common variables between datasets. Can they be used for
  merging? Why or why not?
\item
  Since ISO codes are standardized, download and use the
  \href{https://www.cia.gov/the-world-factbook/references/country-data-codes/}{CIA
  country data codes} for merging. Make sure you are not losing any of
  the countries in your original data set when merging.
\end{itemize}

\begin{Shaded}
\begin{Highlighting}[]
\NormalTok{country\_codes }\OtherTok{\textless{}{-}} \FunctionTok{read\_csv}\NormalTok{(}\StringTok{"Country Data Codes.csv"}\NormalTok{, }\AttributeTok{show\_col\_types =} \ConstantTok{FALSE}\NormalTok{) }\SpecialCharTok{\%\textgreater{}\%}
  \FunctionTok{rename}\NormalTok{(}
    \AttributeTok{country =}\NormalTok{ Name,}
    \AttributeTok{iso\_code =}\NormalTok{ GENC}
\NormalTok{  ) }\SpecialCharTok{\%\textgreater{}\%}
  \FunctionTok{select}\NormalTok{(country, iso\_code) }\SpecialCharTok{\%\textgreater{}\%}
  \FunctionTok{mutate}\NormalTok{(}\AttributeTok{country =} \FunctionTok{str\_trim}\NormalTok{(country))}

\NormalTok{country\_data\_codes }\OtherTok{\textless{}{-}}\NormalTok{ country\_data }\SpecialCharTok{\%\textgreater{}\%}
  \FunctionTok{left\_join}\NormalTok{(country\_codes, }\AttributeTok{by =} \StringTok{"country"}\NormalTok{)}
\FunctionTok{glimpse}\NormalTok{(country\_codes)}
\end{Highlighting}
\end{Shaded}

\begin{verbatim}
## Rows: 278
## Columns: 2
## $ country  <chr> "Afghanistan", "Akrotiri", "Albania", "Algeria", "American Sa~
## $ iso_code <chr> "AFG", "XQZ", "ALB", "DZA", "ASM", "AND", "AGO", "AIA", "ATA"~
\end{verbatim}

\begin{Shaded}
\begin{Highlighting}[]
\FunctionTok{head}\NormalTok{(country\_codes)}
\end{Highlighting}
\end{Shaded}

\begin{verbatim}
## # A tibble: 6 x 2
##   country        iso_code
##   <chr>          <chr>   
## 1 Afghanistan    AFG     
## 2 Akrotiri       XQZ     
## 3 Albania        ALB     
## 4 Algeria        DZA     
## 5 American Samoa ASM     
## 6 Andorra        AND
\end{verbatim}

\begin{Shaded}
\begin{Highlighting}[]
\NormalTok{income\_url }\OtherTok{\textless{}{-}} \StringTok{"https://ourworldindata.org/grapher/world{-}bank{-}income{-}groups.csv?v=1\&csvType=full\&useColumnShortNames=false"}

\NormalTok{income\_classification }\OtherTok{\textless{}{-}} \FunctionTok{read\_csv}\NormalTok{(income\_url, }\AttributeTok{show\_col\_types =} \ConstantTok{FALSE}\NormalTok{) }\SpecialCharTok{\%\textgreater{}\%}
  \FunctionTok{filter}\NormalTok{(Year }\SpecialCharTok{==} \DecValTok{2024}\NormalTok{) }\SpecialCharTok{\%\textgreater{}\%}
  \FunctionTok{rename}\NormalTok{(}\AttributeTok{iso\_code =}\NormalTok{ Code, }
         \AttributeTok{income\_status =} \StringTok{"World Bank\textquotesingle{}s income classification"}\NormalTok{) }\SpecialCharTok{\%\textgreater{}\%}
  \FunctionTok{select}\NormalTok{(iso\_code, income\_status)}
\FunctionTok{glimpse}\NormalTok{(income\_classification)}
\end{Highlighting}
\end{Shaded}

\begin{verbatim}
## Rows: 218
## Columns: 2
## $ iso_code      <chr> "AFG", "ALB", "DZA", "ASM", "AND", "AGO", "ATG", "ARG", ~
## $ income_status <chr> "Low-income countries", "Upper-middle-income countries",~
\end{verbatim}

\begin{Shaded}
\begin{Highlighting}[]
\FunctionTok{head}\NormalTok{(income\_classification)}
\end{Highlighting}
\end{Shaded}

\begin{verbatim}
## # A tibble: 6 x 2
##   iso_code income_status                
##   <chr>    <chr>                        
## 1 AFG      Low-income countries         
## 2 ALB      Upper-middle-income countries
## 3 DZA      Upper-middle-income countries
## 4 ASM      High-income countries        
## 5 AND      High-income countries        
## 6 AGO      Lower-middle-income countries
\end{verbatim}

\begin{Shaded}
\begin{Highlighting}[]
\NormalTok{country\_data\_income }\OtherTok{\textless{}{-}}\NormalTok{ country\_data\_codes }\SpecialCharTok{\%\textgreater{}\%}
  \FunctionTok{left\_join}\NormalTok{(income\_classification, }\AttributeTok{by =} \StringTok{"iso\_code"}\NormalTok{)}
\FunctionTok{glimpse}\NormalTok{(country\_data\_income)}
\end{Highlighting}
\end{Shaded}

\begin{verbatim}
## Rows: 227
## Columns: 6
## $ country                 <chr> "Cuba", "Micronesia, Federated States of", "So~
## $ education_expenditure   <dbl> 12.8, 12.5, 9.9, 8.8, 7.9, 7.8, 7.7, 7.7, 7.6,~
## $ youth_unemployment_rate <dbl> 6.1, 18.9, 1.3, NA, 9.7, 9.4, 6.1, 9.4, 15.3, ~
## $ net_migration_rate      <dbl> -3.7, -20.9, -1.6, 0.0, 4.0, 2.8, 3.3, -1.2, -~
## $ iso_code                <chr> "CUB", "FSM", "SLB", "MSR", "NOR", "DNK", "ISL~
## $ income_status           <chr> "Upper-middle-income countries", "Lower-middle~
\end{verbatim}

\begin{Shaded}
\begin{Highlighting}[]
\FunctionTok{head}\NormalTok{(country\_data\_income)}
\end{Highlighting}
\end{Shaded}

\begin{verbatim}
## # A tibble: 6 x 6
##   country        education_expenditure youth_unemployment_r~1 net_migration_rate
##   <chr>                          <dbl>                  <dbl>              <dbl>
## 1 Cuba                            12.8                    6.1               -3.7
## 2 Micronesia, F~                  12.5                   18.9              -20.9
## 3 Solomon Islan~                   9.9                    1.3               -1.6
## 4 Montserrat                       8.8                   NA                  0  
## 5 Norway                           7.9                    9.7                4  
## 6 Denmark                          7.8                    9.4                2.8
## # i abbreviated name: 1: youth_unemployment_rate
## # i 2 more variables: iso_code <chr>, income_status <chr>
\end{verbatim}

\begin{Shaded}
\begin{Highlighting}[]
\FunctionTok{nrow}\NormalTok{(country\_data) }\SpecialCharTok{==} \FunctionTok{nrow}\NormalTok{(country\_data\_income)}
\end{Highlighting}
\end{Shaded}

\begin{verbatim}
## [1] TRUE
\end{verbatim}

\textbf{Conclusion:}

We can obtain country income classification from the World Bank by
directly calling read\_csv() with the url:
\url{https://ourworldindata.org/grapher/world-bank-income-groups.csv?v=1&csvType=full&useColumnShortNames=false}
. By doing so the .csv is directly loaded into R.

Identify common variables between datasets. Can they be used for
merging? Why or why not?

By inspecting the ``Country Data Codes.csv'' and the
``world-bank-income-group.csv'' loaded via the URL, we see that the GENC
column and Code column are fitting for merging, due to the ISO/GENC
country codes being standardized, therefore the risk of deviating
variables is quite low. Country Name and Entity might appear similar but
should not be used for merging, due to differences in pronunciation,
spelling and naming conventions regarding country names.

Since ISO codes are standardized, download and use the
\href{https://www.cia.gov/the-world-factbook/references/country-data-codes/}{CIA
country data codes} for merging. Make sure you are not losing any of the
countries in your original data set when merging.

By comparing the number of rows of the country\_data and the via
left\_join enriched country\_data\_income tibbles we can check if any
country rows were lost by joining. TRUE is returned, therefore both
tibbles have the same number of rows. country\_data\_income now just
contains the additional income classification information.

\subsection{d.~Adding Geographical
Information}\label{d.-adding-geographical-information}

Introduce continent and subcontinent (or region) data for each country.

\begin{itemize}
\tightlist
\item
  Find and download an appropriate online resource.
\item
  Merge this information into the dataset, naming the final dataset
  \texttt{df\_vars}. Make sure you are not losing any of the countries
  in your original data set when merging.
\end{itemize}

\begin{Shaded}
\begin{Highlighting}[]
\NormalTok{region\_subcontinent\_url }\OtherTok{\textless{}{-}} \StringTok{"https://ourworldindata.org/grapher/continents{-}according{-}to{-}our{-}world{-}in{-}data.csv?v=1\&csvType=full\&useColumnShortNames=false"}

\NormalTok{region\_data }\OtherTok{\textless{}{-}} \FunctionTok{read\_csv}\NormalTok{(region\_subcontinent\_url, }\AttributeTok{show\_col\_types =} \ConstantTok{FALSE}\NormalTok{) }\SpecialCharTok{\%\textgreater{}\%}
    \FunctionTok{rename}\NormalTok{(}
    \AttributeTok{country =}\NormalTok{ Entity,}
    \AttributeTok{iso\_code =}\NormalTok{ Code,}
    \AttributeTok{region =} \StringTok{"World region according to OWID"}
\NormalTok{  ) }\SpecialCharTok{\%\textgreater{}\%}
  \FunctionTok{select}\NormalTok{(iso\_code, region) }\SpecialCharTok{\%\textgreater{}\%}
  \FunctionTok{distinct}\NormalTok{(iso\_code, }\AttributeTok{.keep\_all =} \ConstantTok{TRUE}\NormalTok{) }\CommentTok{\# distinct due to iso\_code finding multiple fitting rows (many{-}to{-}many)}
\FunctionTok{glimpse}\NormalTok{(region\_data)}
\end{Highlighting}
\end{Shaded}

\begin{verbatim}
## Rows: 258
## Columns: 2
## $ iso_code <chr> "AFG", "ALA", "ALB", "DZA", "ASM", "AND", "AGO", "AIA", "ATG"~
## $ region   <chr> "Asia", "Europe", "Europe", "Africa", "Oceania", "Europe", "A~
\end{verbatim}

\begin{Shaded}
\begin{Highlighting}[]
\FunctionTok{head}\NormalTok{(region\_data)}
\end{Highlighting}
\end{Shaded}

\begin{verbatim}
## # A tibble: 6 x 2
##   iso_code region 
##   <chr>    <chr>  
## 1 AFG      Asia   
## 2 ALA      Europe 
## 3 ALB      Europe 
## 4 DZA      Africa 
## 5 ASM      Oceania
## 6 AND      Europe
\end{verbatim}

\begin{Shaded}
\begin{Highlighting}[]
\NormalTok{df\_vars }\OtherTok{\textless{}{-}}\NormalTok{ country\_data\_income }\SpecialCharTok{\%\textgreater{}\%}
  \FunctionTok{left\_join}\NormalTok{(region\_data, }\AttributeTok{by =} \StringTok{"iso\_code"}\NormalTok{)}
\FunctionTok{glimpse}\NormalTok{(df\_vars)}
\end{Highlighting}
\end{Shaded}

\begin{verbatim}
## Rows: 227
## Columns: 7
## $ country                 <chr> "Cuba", "Micronesia, Federated States of", "So~
## $ education_expenditure   <dbl> 12.8, 12.5, 9.9, 8.8, 7.9, 7.8, 7.7, 7.7, 7.6,~
## $ youth_unemployment_rate <dbl> 6.1, 18.9, 1.3, NA, 9.7, 9.4, 6.1, 9.4, 15.3, ~
## $ net_migration_rate      <dbl> -3.7, -20.9, -1.6, 0.0, 4.0, 2.8, 3.3, -1.2, -~
## $ iso_code                <chr> "CUB", "FSM", "SLB", "MSR", "NOR", "DNK", "ISL~
## $ income_status           <chr> "Upper-middle-income countries", "Lower-middle~
## $ region                  <chr> "North America", "Oceania", "Oceania", "North ~
\end{verbatim}

\begin{Shaded}
\begin{Highlighting}[]
\FunctionTok{head}\NormalTok{(df\_vars)}
\end{Highlighting}
\end{Shaded}

\begin{verbatim}
## # A tibble: 6 x 7
##   country        education_expenditure youth_unemployment_r~1 net_migration_rate
##   <chr>                          <dbl>                  <dbl>              <dbl>
## 1 Cuba                            12.8                    6.1               -3.7
## 2 Micronesia, F~                  12.5                   18.9              -20.9
## 3 Solomon Islan~                   9.9                    1.3               -1.6
## 4 Montserrat                       8.8                   NA                  0  
## 5 Norway                           7.9                    9.7                4  
## 6 Denmark                          7.8                    9.4                2.8
## # i abbreviated name: 1: youth_unemployment_rate
## # i 3 more variables: iso_code <chr>, income_status <chr>, region <chr>
\end{verbatim}

\begin{Shaded}
\begin{Highlighting}[]
\FunctionTok{nrow}\NormalTok{(country\_data\_income) }\SpecialCharTok{==} \FunctionTok{nrow}\NormalTok{(df\_vars)}
\end{Highlighting}
\end{Shaded}

\begin{verbatim}
## [1] TRUE
\end{verbatim}

\textbf{Conclusion:}

We can again use the website ``\url{https://ourworldindata.org}'' to
find and download appropriate region and subcontinent data via url and
the read\_csv() function.

The distinct() function had to be used due to left\_join() resulting in
some iso\_codes having multiple fitting rows in the
region\_subcontinent\_data. This would result in a many-to-many
relationship. By applying distinct() only unique rows are joined. Now
country\_data\_income and df\_vars have the same number of rows as
demonstrated by TRUE.

\subsection{e. Data Tidiness and Summary
Statistics}\label{e.-data-tidiness-and-summary-statistics}

\begin{itemize}
\tightlist
\item
  Evaluate the tidiness of \texttt{df\_vars} (observational units,
  variables, fixed vs.~measured variables). Make adjustments to tidy the
  data, if necessary.
\end{itemize}

A dataset is called tidy if following four criteria hold: 1. Each
variable forms a column. 2. Each observation forms a row. 3. Each type
of observational unit forms a table

Furthermore Fixed variables should come before measured variables.

We see in subtask d that the first three criteria are met. The
observational units in our case are the countries and each collumn is a
seperate variable. We should however rearrange the data a bit.

\begin{Shaded}
\begin{Highlighting}[]
\NormalTok{df\_vars }\OtherTok{\textless{}{-}}\NormalTok{ df\_vars }\SpecialCharTok{\%\textgreater{}\%}
  \FunctionTok{select}\NormalTok{(iso\_code, country, region, income\_status, }\FunctionTok{everything}\NormalTok{())}
\FunctionTok{head}\NormalTok{(df\_vars)}
\end{Highlighting}
\end{Shaded}

\begin{verbatim}
## # A tibble: 6 x 7
##   iso_code country                    region income_status education_expenditure
##   <chr>    <chr>                      <chr>  <chr>                         <dbl>
## 1 CUB      Cuba                       North~ Upper-middle~                  12.8
## 2 FSM      Micronesia, Federated Sta~ Ocean~ Lower-middle~                  12.5
## 3 SLB      Solomon Islands            Ocean~ Lower-middle~                   9.9
## 4 MSR      Montserrat                 North~ <NA>                            8.8
## 5 NOR      Norway                     Europe High-income ~                   7.9
## 6 DNK      Denmark                    Europe High-income ~                   7.8
## # i 2 more variables: youth_unemployment_rate <dbl>, net_migration_rate <dbl>
\end{verbatim}

\begin{itemize}
\tightlist
\item
  Create a frequency table for the income status variable and briefly
  interpret the results.
\end{itemize}

\begin{Shaded}
\begin{Highlighting}[]
\NormalTok{df\_vars }\SpecialCharTok{\%\textgreater{}\%}
  \FunctionTok{group\_by}\NormalTok{(income\_status) }\SpecialCharTok{\%\textgreater{}\%}
  \FunctionTok{count}\NormalTok{()}
\end{Highlighting}
\end{Shaded}

\begin{verbatim}
## # A tibble: 5 x 2
## # Groups:   income_status [5]
##   income_status                     n
##   <chr>                         <int>
## 1 High-income countries            86
## 2 Low-income countries             25
## 3 Lower-middle-income countries    49
## 4 Upper-middle-income countries    51
## 5 <NA>                             16
\end{verbatim}

we see that high income countries are most common, followed by
lower/upper middle income countries and at last low income countries.
There are some without classification, which could be explained by the
remotness and difficulty in assesing this in very small
countries/islands.

\begin{itemize}
\tightlist
\item
  Analyze the distribution of income status across continents by
  computing absolute and relative frequencies. Comment on the findings.
\end{itemize}

\begin{Shaded}
\begin{Highlighting}[]
\NormalTok{df\_vars }\SpecialCharTok{\%\textgreater{}\%}
  \FunctionTok{group\_by}\NormalTok{(region, income\_status) }\SpecialCharTok{\%\textgreater{}\%}
  \FunctionTok{summarise}\NormalTok{(}\AttributeTok{abs\_freq =} \FunctionTok{n}\NormalTok{(), }\AttributeTok{rel\_freq =} \FunctionTok{n}\NormalTok{()}\SpecialCharTok{/}\FunctionTok{nrow}\NormalTok{(df\_vars) ) }\SpecialCharTok{\%\textgreater{}\%}
  \FunctionTok{arrange}\NormalTok{(}\FunctionTok{desc}\NormalTok{(rel\_freq)) }\SpecialCharTok{\%\textgreater{}\%}
  \FunctionTok{head}\NormalTok{(}\DecValTok{10}\NormalTok{)}
\end{Highlighting}
\end{Shaded}

\begin{verbatim}
## `summarise()` has regrouped the output.
## i Summaries were computed grouped by region and income_status.
## i Output is grouped by region.
## i Use `summarise(.groups = "drop_last")` to silence this message.
## i Use `summarise(.by = c(region, income_status))` for per-operation grouping
##   (`?dplyr::dplyr_by`) instead.
\end{verbatim}

\begin{verbatim}
## # A tibble: 10 x 4
## # Groups:   region [5]
##    region        income_status                 abs_freq rel_freq
##    <chr>         <chr>                            <int>    <dbl>
##  1 Europe        High-income countries               39   0.172 
##  2 Africa        Lower-middle-income countries       23   0.101 
##  3 Africa        Low-income countries                21   0.0925
##  4 North America High-income countries               20   0.0881
##  5 Asia          Lower-middle-income countries       17   0.0749
##  6 Asia          High-income countries               14   0.0617
##  7 Asia          Upper-middle-income countries       13   0.0573
##  8 North America Upper-middle-income countries       11   0.0485
##  9 Oceania       High-income countries                9   0.0396
## 10 Africa        Upper-middle-income countries        8   0.0352
\end{verbatim}

We observe that most high income countries are in Europe, most low and
lower-middle income countries are in Africa and most upper-middle income
countries are in Asia

\begin{itemize}
\tightlist
\item
  Using the distribution of income status across continents, identify
  which countries are the only ones in their income group across the
  continent. Discuss briefly.
\end{itemize}

\begin{Shaded}
\begin{Highlighting}[]
\NormalTok{df\_vars }\SpecialCharTok{\%\textgreater{}\%} 
  \FunctionTok{group\_by}\NormalTok{(region, income\_status) }\SpecialCharTok{\%\textgreater{}\%}
  \FunctionTok{filter}\NormalTok{(}\FunctionTok{n}\NormalTok{() }\SpecialCharTok{==} \DecValTok{1}\NormalTok{) }\SpecialCharTok{\%\textgreater{}\%}
  \FunctionTok{select}\NormalTok{(country, region, income\_status)}
\end{Highlighting}
\end{Shaded}

\begin{verbatim}
## # A tibble: 3 x 3
## # Groups:   region, income_status [3]
##   country    region        income_status                
##   <chr>      <chr>         <chr>                        
## 1 Bolivia    South America Lower-middle-income countries
## 2 Seychelles Africa        High-income countries        
## 3 Venezuela  South America <NA>
\end{verbatim}

We see that Bolivia is the only lower-middle-income country in South
America. Moreover the Seychelles is somewhat surprisingly the only
High-income country in Africa, which could be because of their role as a
tax haven. Lastly, Venezuela is the only country in South America
without classification, which could be caused by the tense political and
economical situation there.

\subsection{f.~Further Summary Statistics and
Insights}\label{f.-further-summary-statistics-and-insights}

\begin{itemize}
\tightlist
\item
  Create a table of average (mean and median) values for expenditure,
  youth unemployment rate and net migration rate separated into income
  status. Make sure that in the output, the ordering of the income
  classes is proper (i.e., L, LM, UM, H or the other way around).
  Briefly comment the results and any differences between the mean and
  median.
\end{itemize}

\begin{Shaded}
\begin{Highlighting}[]
\NormalTok{df\_vars }\OtherTok{\textless{}{-}}\NormalTok{ df\_vars }\SpecialCharTok{\%\textgreater{}\%}
  \FunctionTok{mutate}\NormalTok{(}\AttributeTok{income\_status =} \FunctionTok{factor}\NormalTok{(income\_status, }
                                \AttributeTok{levels =} \FunctionTok{c}\NormalTok{(}\StringTok{"Low{-}income countries"}\NormalTok{, }
                                           \StringTok{"Lower{-}middle{-}income countries"}\NormalTok{, }
                                           \StringTok{"Upper{-}middle{-}income countries"}\NormalTok{, }
                                           \StringTok{"High{-}income countries"}\NormalTok{)))}

\NormalTok{df\_vars }\SpecialCharTok{\%\textgreater{}\%}
  \FunctionTok{filter}\NormalTok{(}\SpecialCharTok{!}\FunctionTok{is.na}\NormalTok{(income\_status)) }\SpecialCharTok{\%\textgreater{}\%}
  \FunctionTok{group\_by}\NormalTok{(income\_status) }\SpecialCharTok{\%\textgreater{}\%}
  \FunctionTok{summarise}\NormalTok{(}\FunctionTok{across}\NormalTok{(}\FunctionTok{c}\NormalTok{(education\_expenditure, youth\_unemployment\_rate, net\_migration\_rate),}
                   \FunctionTok{list}\NormalTok{(}\AttributeTok{mean =} \SpecialCharTok{\textasciitilde{}}\FunctionTok{mean}\NormalTok{(.x, }\AttributeTok{na.rm =} \ConstantTok{TRUE}\NormalTok{), }
                        \AttributeTok{median =} \SpecialCharTok{\textasciitilde{}}\FunctionTok{median}\NormalTok{(.x, }\AttributeTok{na.rm =} \ConstantTok{TRUE}\NormalTok{))))}
\end{Highlighting}
\end{Shaded}

\begin{verbatim}
## # A tibble: 4 x 7
##   income_status                 education_expenditure_m~1 education_expenditur~2
##   <fct>                                             <dbl>                  <dbl>
## 1 Low-income countries                               3.53                    3.1
## 2 Lower-middle-income countries                      4.43                    3.9
## 3 Upper-middle-income countries                      4.52                    4.1
## 4 High-income countries                              4.6                     4.7
## # i abbreviated names: 1: education_expenditure_mean,
## #   2: education_expenditure_median
## # i 4 more variables: youth_unemployment_rate_mean <dbl>,
## #   youth_unemployment_rate_median <dbl>, net_migration_rate_mean <dbl>,
## #   net_migration_rate_median <dbl>
\end{verbatim}

The table shows that the mean youth unemployment rate is higher than the
median, suggesting a right skewed distribution, coming from the fact
that a few countries with extremly high unemployment rate pull the
average up. Net migration also shows a clear trend, high income
countries tend to have positive net migration (influx of people) while
low income countries have a negative net migration in general (people
leaving). The differences between mean and median again highlight the
presence of extreme outliers.

\begin{itemize}
\tightlist
\item
  Look at the standard deviation and the interquartile range of the
  variables per income status instead of the location statistics above.
  Do you gain additional insights? Briefly comment the results.
\end{itemize}

\begin{Shaded}
\begin{Highlighting}[]
\NormalTok{df\_vars }\SpecialCharTok{\%\textgreater{}\%}
  \FunctionTok{filter}\NormalTok{(}\SpecialCharTok{!}\FunctionTok{is.na}\NormalTok{(income\_status)) }\SpecialCharTok{\%\textgreater{}\%}
  \FunctionTok{group\_by}\NormalTok{(income\_status) }\SpecialCharTok{\%\textgreater{}\%}
  \FunctionTok{summarise}\NormalTok{(}\FunctionTok{across}\NormalTok{(}\FunctionTok{c}\NormalTok{(education\_expenditure, youth\_unemployment\_rate, net\_migration\_rate),}
                   \FunctionTok{list}\NormalTok{(}\AttributeTok{stand\_dev =} \SpecialCharTok{\textasciitilde{}}\FunctionTok{sd}\NormalTok{(.x, }\AttributeTok{na.rm =} \ConstantTok{TRUE}\NormalTok{),}
                        \AttributeTok{iqr =} \SpecialCharTok{\textasciitilde{}}\FunctionTok{IQR}\NormalTok{(.x, }\AttributeTok{na.rm =} \ConstantTok{TRUE}\NormalTok{))))}
\end{Highlighting}
\end{Shaded}

\begin{verbatim}
## # A tibble: 4 x 7
##   income_status                 education_expenditure_s~1 education_expenditur~2
##   <fct>                                             <dbl>                  <dbl>
## 1 Low-income countries                               1.73                   2.65
## 2 Lower-middle-income countries                      2.24                   2.55
## 3 Upper-middle-income countries                      1.80                   1.68
## 4 High-income countries                              1.49                   1.8 
## # i abbreviated names: 1: education_expenditure_stand_dev,
## #   2: education_expenditure_iqr
## # i 4 more variables: youth_unemployment_rate_stand_dev <dbl>,
## #   youth_unemployment_rate_iqr <dbl>, net_migration_rate_stand_dev <dbl>,
## #   net_migration_rate_iqr <dbl>
\end{verbatim}

We see that youth unemployment varies substantially across all groups,
meaning income status alone does not guarantee a stable youth labor
market. Similarly the high migration rate standard deviation indicates
differing behavior from countries in high income groups, despite the
mean being high in this group.

\begin{itemize}
\tightlist
\item
  Extend the analysis of the statistics median and IQR to \textbf{each
  income status and continent combination}. Play around with displaying
  the resulting table. Use \texttt{pivot\_longer()} and/or
  \texttt{pivot\_wider()} to generate different outputs. Discuss the
  results as well as the readability of the different tables.
\end{itemize}

\begin{Shaded}
\begin{Highlighting}[]
\NormalTok{stat\_table }\OtherTok{\textless{}{-}}\NormalTok{ df\_vars }\SpecialCharTok{\%\textgreater{}\%}
  \FunctionTok{filter}\NormalTok{(}\SpecialCharTok{!}\FunctionTok{is.na}\NormalTok{(income\_status) }\SpecialCharTok{\&} \SpecialCharTok{!}\FunctionTok{is.na}\NormalTok{(region)) }\SpecialCharTok{\%\textgreater{}\%}
  \FunctionTok{group\_by}\NormalTok{(region, income\_status) }\SpecialCharTok{\%\textgreater{}\%}
  \FunctionTok{summarise}\NormalTok{(}\FunctionTok{across}\NormalTok{(}\FunctionTok{c}\NormalTok{(education\_expenditure, youth\_unemployment\_rate, net\_migration\_rate),}
                   \FunctionTok{list}\NormalTok{(}\AttributeTok{median =} \SpecialCharTok{\textasciitilde{}}\FunctionTok{median}\NormalTok{(.x, }\AttributeTok{na.rm =} \ConstantTok{TRUE}\NormalTok{), }
                        \AttributeTok{iqr =} \SpecialCharTok{\textasciitilde{}}\FunctionTok{IQR}\NormalTok{(.x, }\AttributeTok{na.rm =} \ConstantTok{TRUE}\NormalTok{))))}
\end{Highlighting}
\end{Shaded}

\begin{verbatim}
## `summarise()` has regrouped the output.
## i Summaries were computed grouped by region and income_status.
## i Output is grouped by region.
## i Use `summarise(.groups = "drop_last")` to silence this message.
## i Use `summarise(.by = c(region, income_status))` for per-operation grouping
##   (`?dplyr::dplyr_by`) instead.
\end{verbatim}

\begin{Shaded}
\begin{Highlighting}[]
\FunctionTok{head}\NormalTok{(stat\_table)}
\end{Highlighting}
\end{Shaded}

\begin{verbatim}
## # A tibble: 6 x 8
## # Groups:   region [2]
##   region income_status             education_expenditur~1 education_expenditur~2
##   <chr>  <fct>                                      <dbl>                  <dbl>
## 1 Africa Low-income countries                        2.95                   2.82
## 2 Africa Lower-middle-income coun~                   3.55                   2.05
## 3 Africa Upper-middle-income coun~                   4.95                   1.33
## 4 Africa High-income countries                       4.4                    0   
## 5 Asia   Low-income countries                        4.1                    0   
## 6 Asia   Lower-middle-income coun~                   3.45                   2.8 
## # i abbreviated names: 1: education_expenditure_median,
## #   2: education_expenditure_iqr
## # i 4 more variables: youth_unemployment_rate_median <dbl>,
## #   youth_unemployment_rate_iqr <dbl>, net_migration_rate_median <dbl>,
## #   net_migration_rate_iqr <dbl>
\end{verbatim}

\begin{Shaded}
\begin{Highlighting}[]
\NormalTok{long\_table }\OtherTok{\textless{}{-}}\NormalTok{ stat\_table }\SpecialCharTok{\%\textgreater{}\%}
  \FunctionTok{pivot\_longer}\NormalTok{(}\SpecialCharTok{{-}}\FunctionTok{c}\NormalTok{(region, income\_status), }\AttributeTok{values\_to =} \StringTok{"value"}\NormalTok{, }\AttributeTok{names\_to =} \StringTok{"metric"}\NormalTok{)}
\FunctionTok{head}\NormalTok{(long\_table)}
\end{Highlighting}
\end{Shaded}

\begin{verbatim}
## # A tibble: 6 x 4
## # Groups:   region [1]
##   region income_status        metric                         value
##   <chr>  <fct>                <chr>                          <dbl>
## 1 Africa Low-income countries education_expenditure_median    2.95
## 2 Africa Low-income countries education_expenditure_iqr       2.82
## 3 Africa Low-income countries youth_unemployment_rate_median  8.7 
## 4 Africa Low-income countries youth_unemployment_rate_iqr    12.4 
## 5 Africa Low-income countries net_migration_rate_median      -0.9 
## 6 Africa Low-income countries net_migration_rate_iqr          3.2
\end{verbatim}

\begin{Shaded}
\begin{Highlighting}[]
\NormalTok{wide\_table }\OtherTok{\textless{}{-}}\NormalTok{ long\_table }\SpecialCharTok{\%\textgreater{}\%}
  \FunctionTok{pivot\_wider}\NormalTok{( }\AttributeTok{names\_from =} \StringTok{"income\_status"}\NormalTok{, }\AttributeTok{values\_from =} \StringTok{"value"}\NormalTok{)}
\FunctionTok{head}\NormalTok{(wide\_table)}
\end{Highlighting}
\end{Shaded}

\begin{verbatim}
## # A tibble: 6 x 6
## # Groups:   region [1]
##   region metric                    `Low-income countries` Lower-middle-income ~1
##   <chr>  <chr>                                      <dbl>                  <dbl>
## 1 Africa education_expenditure_me~                   2.95                   3.55
## 2 Africa education_expenditure_iqr                   2.82                   2.05
## 3 Africa youth_unemployment_rate_~                   8.7                   16.5 
## 4 Africa youth_unemployment_rate_~                  12.4                   22.2 
## 5 Africa net_migration_rate_median                  -0.9                   -0.4 
## 6 Africa net_migration_rate_iqr                      3.2                    1.65
## # i abbreviated name: 1: `Lower-middle-income countries`
## # i 2 more variables: `Upper-middle-income countries` <dbl>,
## #   `High-income countries` <dbl>
\end{verbatim}

The standard summarise output is a wide table with many columns which
could be difficult to read. We transform it first to a long format,
which is less intuitive for humans to read but more suitable for further
analysis and plotting. Lastly we transform it to another wide table,
which is much easier to read for comparative analysis.

\begin{itemize}
\tightlist
\item
  Identify countries performing well in terms of both \textbf{youth
  unemployment} and \textbf{net migration rate} (top 25\% in net
  migration and bottom 25\% in youth unemployment within their
  continent).
\end{itemize}

\begin{Shaded}
\begin{Highlighting}[]
\NormalTok{df\_vars }\SpecialCharTok{\%\textgreater{}\%}
  \FunctionTok{group\_by}\NormalTok{(region) }\SpecialCharTok{\%\textgreater{}\%}
  \FunctionTok{filter}\NormalTok{(net\_migration\_rate }\SpecialCharTok{\textgreater{}=} \FunctionTok{quantile}\NormalTok{(net\_migration\_rate, }\FloatTok{0.75}\NormalTok{, }\AttributeTok{na.rm =} \ConstantTok{TRUE}\NormalTok{) }\SpecialCharTok{\&}
\NormalTok{         youth\_unemployment\_rate }\SpecialCharTok{\textless{}=} \FunctionTok{quantile}\NormalTok{(youth\_unemployment\_rate, }\FloatTok{0.25}\NormalTok{, }\AttributeTok{na.rm =} \ConstantTok{TRUE}\NormalTok{)) }\SpecialCharTok{\%\textgreater{}\%}
  \FunctionTok{select}\NormalTok{(region, country, youth\_unemployment\_rate, net\_migration\_rate) }\SpecialCharTok{\%\textgreater{}\%}
  \FunctionTok{arrange}\NormalTok{(region, youth\_unemployment\_rate)}
\end{Highlighting}
\end{Shaded}

\begin{verbatim}
## # A tibble: 18 x 4
## # Groups:   region [6]
##    region        country              youth_unemployment_rate net_migration_rate
##    <chr>         <chr>                                  <dbl>              <dbl>
##  1 Africa        Madagascar                               1                  0  
##  2 Africa        Guinea                                   1                  0  
##  3 Africa        Togo                                     3.9                0  
##  4 Africa        Cote d'Ivoire                            5.5                1.2
##  5 Africa        Benin                                    5.6                0.3
##  6 Asia          Qatar                                    0.4                6.5
##  7 Asia          Kazakhstan                               3.8                0.4
##  8 Asia          Macau                                    5.3                3.3
##  9 Asia          Bahrain                                  5.3               10.6
## 10 Asia          United Arab Emirates                     6.9                7.6
## 11 Europe        Switzerland                              7.9                4.6
## 12 Europe        Malta                                    9.1                6.6
## 13 Europe        Norway                                   9.7                4  
## 14 North America United States                            8.6                3  
## 15 Oceania       Papua New Guinea                         3.6                0  
## 16 Oceania       Palau                                    5.6                0.9
## 17 South America Ecuador                                  7.9                0  
## 18 South America Suriname                                13.4                0.5
\end{verbatim}

\subsection{g. Conditional
Probabilities}\label{g.-conditional-probabilities}

Estimate the following based on the observed frequencies in the data:

\begin{itemize}
\tightlist
\item
  What is the (posterior or conditional) probability that a European
  country belongs to the high income group? What is the prior
  probability that a country belongs to the high income group?
\end{itemize}

\begin{Shaded}
\begin{Highlighting}[]
\CommentTok{\# Probability of a country being High Income given it is European}
\NormalTok{df\_vars }\SpecialCharTok{\%\textgreater{}\%}
  \FunctionTok{filter}\NormalTok{(region }\SpecialCharTok{==} \StringTok{"Europe"}\NormalTok{) }\SpecialCharTok{\%\textgreater{}\%}
   \FunctionTok{summarise}\NormalTok{(}\AttributeTok{Pct\_European\_High\_Income =} \DecValTok{100}\SpecialCharTok{*}\FunctionTok{mean}\NormalTok{(income\_status }\SpecialCharTok{==} \StringTok{"High{-}income countries"}\NormalTok{,}\AttributeTok{na.rm=}\ConstantTok{TRUE}\NormalTok{))}
\end{Highlighting}
\end{Shaded}

\begin{verbatim}
## # A tibble: 1 x 1
##   Pct_European_High_Income
##                      <dbl>
## 1                     84.8
\end{verbatim}

\begin{Shaded}
\begin{Highlighting}[]
\CommentTok{\# Probability of any country being High Income}
\NormalTok{df\_vars }\SpecialCharTok{\%\textgreater{}\%}
   \FunctionTok{summarise}\NormalTok{(}\AttributeTok{Pct\_Global\_High\_Income =} \DecValTok{100}\SpecialCharTok{*}\FunctionTok{mean}\NormalTok{(income\_status }\SpecialCharTok{==} \StringTok{"High{-}income countries"}\NormalTok{,}\AttributeTok{na.rm=}\ConstantTok{TRUE}\NormalTok{))}
\end{Highlighting}
\end{Shaded}

\begin{verbatim}
## # A tibble: 1 x 1
##   Pct_Global_High_Income
##                    <dbl>
## 1                   40.8
\end{verbatim}

We can observe that the conditional probability is more than double the
unconditional. This is unsurprising as many of the worlds highest income
countries are located in Europe.

\begin{itemize}
\tightlist
\item
  Given a country has high youth unemployment (above \%25), what is the
  probability that it also has negative net migration?
\end{itemize}

\begin{Shaded}
\begin{Highlighting}[]
\CommentTok{\# Probability of a country having negative net migration given it has high youth unemployment}
\NormalTok{df\_vars }\SpecialCharTok{\%\textgreater{}\%}
  \FunctionTok{filter}\NormalTok{(youth\_unemployment\_rate }\SpecialCharTok{\textgreater{}} \DecValTok{25}\NormalTok{) }\SpecialCharTok{\%\textgreater{}\%}
   \FunctionTok{summarise}\NormalTok{(}\AttributeTok{Pct\_Negative\_Net\_Migration =} \DecValTok{100}\SpecialCharTok{*}\FunctionTok{mean}\NormalTok{(net\_migration\_rate }\SpecialCharTok{\textless{}} \DecValTok{0}\NormalTok{,}\AttributeTok{na.rm=}\ConstantTok{TRUE}\NormalTok{))}
\end{Highlighting}
\end{Shaded}

\begin{verbatim}
## # A tibble: 1 x 1
##   Pct_Negative_Net_Migration
##                        <dbl>
## 1                       61.2
\end{verbatim}

The conditional probability is around 61\%.

\subsection{h. Simpson's Paradox
Analysis}\label{h.-simpsons-paradox-analysis}

Investigate whether an overall trend in youth unemployment rate in the
high and low income groups reverses when analyzed at the continent
level. E.g., does the youth unemployment rate appear lower in low-income
countries overall, but higher when controlling for continent? Explain
the results and possible reasons behind this paradox.

\begin{Shaded}
\begin{Highlighting}[]
\CommentTok{\# Overall trend for income groups}
\NormalTok{df\_vars }\SpecialCharTok{\%\textgreater{}\%}
 \FunctionTok{drop\_na}\NormalTok{(income\_status) }\SpecialCharTok{\%\textgreater{}\%}
  \FunctionTok{group\_by}\NormalTok{(income\_status) }\SpecialCharTok{\%\textgreater{}\%}
    \FunctionTok{summarise}\NormalTok{(}\AttributeTok{mean\_youth\_unemployment\_rate =} \FunctionTok{mean}\NormalTok{(youth\_unemployment\_rate,}\AttributeTok{na.rm=}\ConstantTok{TRUE}\NormalTok{))}
\end{Highlighting}
\end{Shaded}

\begin{verbatim}
## # A tibble: 4 x 2
##   income_status                 mean_youth_unemployment_rate
##   <fct>                                                <dbl>
## 1 Low-income countries                                  14.9
## 2 Lower-middle-income countries                         15.8
## 3 Upper-middle-income countries                         21.6
## 4 High-income countries                                 16.5
\end{verbatim}

\begin{Shaded}
\begin{Highlighting}[]
\CommentTok{\# Trend for different continents}
\NormalTok{df\_vars }\SpecialCharTok{\%\textgreater{}\%}
  \FunctionTok{drop\_na}\NormalTok{(income\_status) }\SpecialCharTok{\%\textgreater{}\%}
   \FunctionTok{pivot\_wider}\NormalTok{(}\AttributeTok{names\_from =}\NormalTok{ region, }\AttributeTok{values\_from =}\NormalTok{ youth\_unemployment\_rate, }
               \AttributeTok{names\_prefix =} \StringTok{"Youth\_UNEMP\_"}\NormalTok{) }\SpecialCharTok{\%\textgreater{}\%}
    \FunctionTok{group\_by}\NormalTok{(income\_status) }\SpecialCharTok{\%\textgreater{}\%}
      \FunctionTok{summarise}\NormalTok{(}\AttributeTok{AvY\_UN\_North\_America =} \FunctionTok{mean}\NormalTok{(}\StringTok{\textasciigrave{}}\AttributeTok{Youth\_UNEMP\_North America}\StringTok{\textasciigrave{}}\NormalTok{,}\AttributeTok{na.rm=}\ConstantTok{TRUE}\NormalTok{),}
                \AttributeTok{AvY\_UN\_Oceania       =} \FunctionTok{mean}\NormalTok{(Youth\_UNEMP\_Oceania,}\AttributeTok{na.rm=}\ConstantTok{TRUE}\NormalTok{),}
                \AttributeTok{AvY\_UN\_Europe        =} \FunctionTok{mean}\NormalTok{(Youth\_UNEMP\_Europe,}\AttributeTok{na.rm=}\ConstantTok{TRUE}\NormalTok{),}
                \AttributeTok{AvY\_UN\_Africa        =} \FunctionTok{mean}\NormalTok{(Youth\_UNEMP\_Africa,}\AttributeTok{na.rm=}\ConstantTok{TRUE}\NormalTok{),}
                \AttributeTok{AvY\_UN\_South\_America =} \FunctionTok{mean}\NormalTok{(}\StringTok{\textasciigrave{}}\AttributeTok{Youth\_UNEMP\_South America}\StringTok{\textasciigrave{}}\NormalTok{,}\AttributeTok{na.rm=}\ConstantTok{TRUE}\NormalTok{),}
                \AttributeTok{AvY\_UN\_Asia          =} \FunctionTok{mean}\NormalTok{(Youth\_UNEMP\_Asia,}\AttributeTok{na.rm=}\ConstantTok{TRUE}\NormalTok{)) }\SpecialCharTok{\%\textgreater{}\%} \FunctionTok{kbl}\NormalTok{()}
\end{Highlighting}
\end{Shaded}

\begin{tabular}[t]{l|r|r|r|r|r|r}
\hline
income\_status & AvY\_UN\_North\_America & AvY\_UN\_Oceania & AvY\_UN\_Europe & AvY\_UN\_Africa & AvY\_UN\_South\_America & AvY\_UN\_Asia\\
\hline
Low-income countries & NaN & NaN & NaN & 12.63333 & NaN & 25.96667\\
\hline
Lower-middle-income countries & 9.60000 & 10.22500 & NaN & 19.16190 & 6.90000 & 13.86429\\
\hline
Upper-middle-income countries & 15.85000 & 19.72500 & 22.95714 & 37.82857 & 17.31429 & 18.18182\\
\hline
High-income countries & 19.41818 & 25.71429 & 15.50571 & 11.60000 & 21.83333 & 11.03846\\
\hline
\end{tabular}

The overall trend is that low income countries have lower youth
unemployment rates and especially upper-middle-income countries have a
very high youth unemployment rate.

When analy \#\# i. Data Export

Export the final tidy dataset from e. as a \textbf{CSV} with:\texttt{;}
as a separator; \texttt{.} representing missing values; no row names
included. Upload the \texttt{.csv} to TUWEL, together with the
submission.

\end{document}
